\textbf{Phần 1. Trắc nghiệm lựa chọn}
\setcounterpageref{ex}{0}

\begin{ex} %[3D1V60]
Điền dấu thích hợp vào chỗ chấm: $68\,540 \dots 68\,499$.
\choice
{$<$}
{\True $>$}
{$=$}
{Không so sánh được}
\loigiai{
Hàng chục nghìn và hàng nghìn bằng nhau; hàng trăm: $5 > 4$ nên $68\,540 > 68\,499$.
\textbf{Đáp án của câu hỏi là B}
}
\end{ex}

\begin{ex} %[3D1V60]
Trong các số sau, số nào lớn nhất: $78\,450, 78\,540, 87\,450, 87\,504$?
\choice
{$78\,540$}
{$87\,450$}
{$78\,450$}
{\True $87\,504$}
\loigiai{
Hàng chục nghìn là 8; hàng trăm: $5 > 4$ nên $87\,504$ lớn nhất.
\textbf{Đáp án của câu hỏi là D}
}
\end{ex}

\begin{ex} %[3D1V60]
Tìm chữ số $x$ để: $45\,x68 < 45\,168$. Chữ số $x$ là
\choice
{\True $0$}
{$1$}
{$2$}
{Không có}
\loigiai{
Để $45\,x68 < 45\,168$ thì $x < 1 \Rightarrow x = 0$.
\textbf{Đáp án của câu hỏi là A}
}
\end{ex}

\begin{ex} %[3D1V60]
Sắp xếp theo thứ tự từ bé đến lớn: $35\,000; 34\,999; 35\,010; 34\,990$.
\choice
{$34\,999; 34\,990; 35\,000; 35\,010$}
{$35\,010; 35\,000; 34\,999; 34\,990$}
{\True $34\,990; 34\,999; 35\,000; 35\,010$}
{$34\,990; 35\,000; 34\,999; 35\,010$}
\loigiai{
So sánh các hàng: $34\,990 < 34\,999 < 35\,000 < 35\,010$.
\textbf{Đáp án của câu hỏi là C}
}
\end{ex}

\begin{ex} %[3D1V60]
Có bao nhiêu số tự nhiên $y$ có năm chữ số thỏa mãn: $89\,998 < y < 90\,004$?
\choice
{$4$ số}
{\True $5$ số}
{$6$ số}
{$7$ số}
\loigiai{
Các số $y$: $89\,999, 90\,000, 90\,001, 90\,002, 90\,003$ (có 5 số).
\textbf{Đáp án của câu hỏi là B}
}
\end{ex}

\begin{ex} %[3D1V60]
Số bé nhất có năm chữ số khác nhau mà chữ số hàng chục nghìn bằng 6 là
\choice
{$61\,023$}
{$60\,234$}
{$60\,000$}
{\True $60\,123$}
\loigiai{
Hàng chục nghìn là 6, các chữ số tiếp theo lấy nhỏ nhất khác nhau: 0, 1, 2, 3: $60\,123$.
\textbf{Đáp án của câu hỏi là D}
}
\end{ex}

\begin{ex} %[3D1V60]
So sánh $A = 50\,000 + 4000 + 300 + 20$ và $B = 54\,325$. Khẳng định đúng là
\choice
{\True $A < B$}
{$A > B$}
{$A = B$}
{Không so sánh được}
\loigiai{
$A = 54\,320$. Vì $54\,320 < 54\,325$ nên $A < B$.
\textbf{Đáp án của câu hỏi là A}
}
\end{ex}

\begin{ex} %[3D1V60]
Dân số của ba phường: Phường Một có $42\,500$ người, Phường Hai có $42\,050$ người, Phường Ba có $43\,000$ người. Phường có ít dân nhất là
\choice
{Phường Một}
{Phường Ba}
{\True Phường Hai}
{Ba phường bằng nhau}
\loigiai{
$42\,050 < 42\,500 < 43\,000$. Phường Hai có ít dân nhất.
\textbf{Đáp án của câu hỏi là C}
}
\end{ex}

\begin{ex} %[3D1V60]
Tìm chữ số thích hợp điền vào dấu $*$ để: $78\,6*2 > 78\,682$.
\choice
{$8$}
{\True $9$}
{$7$}
{$6$}
\loigiai{
Để $78\,6*2 > 78\,682$ thì $* > 8 \Rightarrow * = 9$.
\textbf{Đáp án của câu hỏi là B}
}
\end{ex}

\begin{ex} %[3D1V60]
Số nào sau đây bé hơn $10\,000$?
\choice
{$10\,001$}
{$10\,000$}
{$10\,010$}
{\True $9999$}
\loigiai{
Số 9999 có 4 chữ số, bé hơn số có 5 chữ số $10\,000$.
\textbf{Đáp án của câu hỏi là D}
}
\end{ex}

\begin{ex} %[3D1V60]
Số lớn nhất có năm chữ số có tổng các chữ số bằng 10 là
\choice
{\True $91\,000$}
{$90\,100$}
{$82\,000$}
{$90\,010$}
\loigiai{
Hàng chục nghìn lớn nhất là 9, hàng nghìn $10 - 9 = 1$, các hàng còn lại là 0: $91\,000$.
\textbf{Đáp án của câu hỏi là A}
}
\end{ex}

\begin{ex} %[3D1V60]
Tìm khoảng giá trị: có bao nhiêu số tròn nghìn $x$ thỏa mãn: $65\,000 < x < 70\,000$?
\choice
{$5$ số}
{$3$ số}
{\True $4$ số}
{$6$ số}
\loigiai{
Các số: $66\,000; 67\,000; 68\,000; 69\,000$ (4 số).
\textbf{Đáp án của câu hỏi là C}
}
\end{ex}

\textbf{Phần 2. Trắc nghiệm đúng sai}
\setcounterpageref{ex}{0}

\begin{ex} %[3D1V60]
Xét tính đúng sai của việc so sánh các số trong phạm vi 100 000:
\choiceTF
{\True Khi so sánh hai số có năm chữ số, ta bắt đầu so sánh từ hàng chục nghìn.}
{Số nào có chữ số hàng đơn vị lớn hơn thì luôn lớn hơn.}
{\True $99\,999 < 100\,000$.}
{$75\,400 > 75\,410$.}
\loigiai{
\begin{itemize}
    \item Ý a) Đúng: Quy tắc so sánh chuẩn.
    \item Ý b) Sai: Sai vì hàng cao nhất mới quyết định.
    \item Ý c) Đúng: Số 5 chữ số luôn bé hơn số 6 chữ số.
    \item Ý d) Sai: Sai vì hàng chục $0 < 1$.
\end{itemize}
\begin{itemize}
    \item a) \textbf{Đúng}.
    \item b) \textbf{Sai}.
    \item c) \textbf{Đúng}.
    \item d) \textbf{Sai}.
\end{itemize}
}
\end{ex}

\begin{ex} %[3D1V60]
Cho bốn số: $64\,800; 64\,080; 68\,400; 60\,840$:
\choiceTF
{\True Số bé nhất trong bốn số là 60 840.}
{\True Số lớn nhất trong bốn số là 68 400.}
{Sắp xếp theo thứ tự tăng dần là: $68\,400, 64\,800, 64\,080, 60\,840$.}
{\True Hiệu giữa số lớn nhất và số bé nhất là: $68\,400 - 60\,840 = 7560$.}
\loigiai{
\begin{itemize}
    \item Ý a) Đúng: Hàng nghìn bé nhất là 0.
    \item Ý b) Đúng: Hàng nghìn lớn nhất là 8.
    \item Ý c) Sai: Sai vì đây là thứ tự giảm dần.
    \item Ý d) Đúng: Phép trừ đúng.
\end{itemize}
\begin{itemize}
    \item a) \textbf{Đúng}.
    \item b) \textbf{Đúng}.
    \item c) \textbf{Sai}.
    \item d) \textbf{Đúng}.
\end{itemize}
}
\end{ex}

\textbf{Phần 3. Trả lời ngắn}
\setcounterpageref{ex}{0}

\begin{ex} %[3D1V60]
Tìm số lớn nhất trong các số sau: $86\,540, 86\,450, 85\,640, 86\,545$.
\shortans{86545}
\loigiai{
\textbf{Đáp số:} $86545$\\
86 545 có chữ số hàng đơn vị là 5 lớn nhất khi các hàng khác bằng nhau.
}
\end{ex}

\begin{ex} %[3D1V60]
Tìm số tự nhiên $x$ thỏa mãn: $49\,999 < x < 50\,001$.
\shortans{50000}
\loigiai{
\textbf{Đáp số:} $50000$\\
$x = 50\,000$.
}
\end{ex}

\begin{ex} %[3D1V60]
Có bao nhiêu số tự nhiên có năm chữ số lớn hơn $99\,990$?
\shortans{9}
\loigiai{
\textbf{Đáp số:} $9$\\
Từ 99 991 đến 99 999 có 9 số.
}
\end{ex}

\begin{ex} %[3D1V60]
Tìm chữ số $a$ lớn nhất để: $54\,a20 < 54\,520$.
\shortans{4}
\loigiai{
\textbf{Đáp số:} $4$\\
Để $a < 5$ thì chữ số lớn nhất là $a = 4$.
}
\end{ex}

\begin{ex} %[3D1V60]
Viết số bé nhất có năm chữ số khác nhau có tổng các chữ số bằng 15.
\shortans{10239}
\loigiai{
\textbf{Đáp số:} $10239$\\
Nghìn 1, trăm 0, chục 2, đơn vị 3, cuối cùng $15 - 6 = 9$: $10\,239$.
}
\end{ex}

\begin{ex} %[3D1V60]
Tính khoảng cách: số $100\,000$ lớn hơn số $95\,000$ bao nhiêu đơn vị?
\shortans{5000}
\loigiai{
\textbf{Đáp số:} $5000$\\
$100\,000 - 95\,000 = 5000$.
}
\end{ex}
